\cleardoublepage
\chapter*{How to use this book}
\addcontentsline{toc}{chapter}{How to use this book}
\markboth{How to use this book}{How to use this book}

Every one of the twenty topics in this book is a self-contained lesson
followed by a full two-module problem set, laid out exactly the way the
Digital SAT presents mathematics.

\section*{The shape of a topic}

\begin{satconcept}[What each topic contains]
\begin{list}{\textbf{\color{satteal}\textbullet}}%
  {\setlength{\leftmargin}{16pt}\setlength{\itemsep}{3pt}}
\item \textbf{The lesson.} The idea, the formulas worth memorising, the traps
      that cost points, and fully worked examples.
\item \textbf{Module 1 --- questions 1 to 22.} The full range of the topic,
      ramping from easy to hard, exactly as the first adaptive module does.
\item \textbf{Module 2 --- questions 23 to 44.} The follow-up module. It
      continues the numbering and pushes into the harder, mixed-skill end of
      the topic.
\item \textbf{Worked solutions.} Every one of the 44 questions is solved in
      full at the back of the book, not merely answered.
\end{list}
\end{satconcept}

\section*{Reading the labels}

Each question carries two small labels. The first is its difficulty
(\satdiffchip{E}, \satdiffchip{M} or \satdiffchip{H}) --- the same three-band
scale the test itself uses to build an adaptive module. The second appears
only on student-produced-response questions
(\sattypechip{SPR}); those have no answer choices and you enter the value
directly. Everything else is four-choice multiple choice.

\section*{How to work through it}

\begin{satstrategy}[A routine that works]
Do Module 1 of a topic untimed and check it against the worked solutions
before moving on. Only once you can finish Module 1 cleanly should you time
Module 2 --- give yourself 22 questions in 35 minutes, which is the real
per-module pace. A question you got right by guessing is a question you got
wrong; mark it and come back.
\end{satstrategy}

\begin{sattrap}[The mistake almost everyone makes]
Reading a worked solution and understanding it is not the same as being able
to produce it. If you looked at the solution, the question is not finished ---
redo it from a blank page a day later.
\end{sattrap}

\section*{What is tested}

The Digital SAT's mathematics section draws on four domains. The twenty
topics of this book cover them in the proportions the test actually uses:

\begin{satformula}[The four domains]
\begin{list}{\textbf{\color{satteal}\textbullet}}%
  {\setlength{\leftmargin}{16pt}\setlength{\itemsep}{3pt}}
\item \textbf{Algebra} --- linear equations, functions, systems and
      inequalities.
\item \textbf{Advanced Math} --- equivalent expressions, nonlinear equations
      and nonlinear functions.
\item \textbf{Problem-Solving and Data Analysis} --- ratios, percentages,
      distributions, scatterplots, probability and inference.
\item \textbf{Geometry and Trigonometry} --- area and volume, lines and
      angles, right triangles and circles.
\end{list}
\end{satformula}
